\documentclass[11pt,letterpaper,oneside]{report}
\usepackage[T1]{fontenc}
\usepackage[utf8]{inputenc}
\usepackage{lmodern}
\usepackage[letterpaper,margin=0.92in,headheight=15pt]{geometry}
\usepackage{amsmath,amssymb,amsthm,mathtools}
\usepackage{microtype}
\usepackage[dvipsnames]{xcolor}
\usepackage{graphicx}
\usepackage{booktabs,longtable,array,tabularx}
\usepackage{enumitem}
\usepackage{fancyhdr}
\usepackage{titlesec}
\usepackage[most]{tcolorbox}
\usepackage{listings}
\usepackage{tikz}
\usetikzlibrary{positioning}
\usepackage{hyperref}
\hypersetup{colorlinks=true,linkcolor=MidnightBlue,urlcolor=RoyalBlue,pdfauthor={Independent Mathematics and Engineering Study Suite},pdftitle={Numerical Methods and Scientific Computing Guide}}
\definecolor{ink}{HTML}{162639}
\definecolor{navy}{HTML}{173B61}
\definecolor{blue}{HTML}{2874A6}
\definecolor{teal}{HTML}{007C83}
\definecolor{amber}{HTML}{8B5E16}
\definecolor{pale}{HTML}{EDF4F8}
\definecolor{paleteal}{HTML}{ECF8F5}
\definecolor{paleamber}{HTML}{FFF6E8}
\titleformat{\chapter}[display]{\normalfont\huge\bfseries\color{navy}}{\filleft\fontsize{43}{46}\selectfont\thechapter}{0.4ex}{\titlerule\vspace{1ex}}[\vspace{1ex}\titlerule]
\titleformat{\section}{\normalfont\Large\bfseries\color{navy}}{\thesection}{0.8em}{}
\titleformat{\subsection}{\normalfont\large\bfseries\color{teal}}{\thesubsection}{0.75em}{}
\pagestyle{fancy}\fancyhf{}\fancyhead[L]{\small\color{navy}Differential Equations Engineering Suite}\fancyhead[R]{\small\color{navy}Volume 4 of 8}\fancyfoot[C]{\thepage}\renewcommand{\headrulewidth}{0.3pt}
\setlist[itemize]{leftmargin=*,itemsep=3pt,topsep=5pt}\setlist[enumerate]{leftmargin=*,itemsep=4pt,topsep=5pt}
\setlength{\parskip}{5pt}\setlength{\parindent}{0pt}\setlength{\emergencystretch}{2em}
\newtheorem{theorem}{Theorem}[chapter]
\newtheorem{proposition}[theorem]{Proposition}
\newtheorem{lemma}[theorem]{Lemma}
\theoremstyle{definition}\newtheorem{definition}[theorem]{Definition}
\newtheorem{example}[theorem]{Worked example}
\theoremstyle{remark}\newtheorem{remark}[theorem]{Remark}
\newtcolorbox{bookground}[1][]{enhanced,breakable,colback=pale,colframe=blue,boxrule=0.6pt,title={TEXTBOOK ALIGNMENT},fonttitle=\sffamily\bfseries\small,#1}
\newtcolorbox{modern}[1][]{enhanced,breakable,colback=paleteal,colframe=teal,boxrule=0.6pt,title={MODERN COMPUTATIONAL EXTENSION},fonttitle=\sffamily\bfseries\small,#1}
\newtcolorbox{caution}[1][]{enhanced,breakable,colback=paleamber,colframe=amber,boxrule=0.6pt,title={ASSUMPTION / VALIDATION},fonttitle=\sffamily\bfseries\small,#1}
\newcommand{\dd}{\,\mathrm{d}}
\newcommand{\R}{\mathbb{R}}
\newcommand{\lap}{\mathcal{L}}
\newcommand{\norm}[1]{\left\lVert#1\right\rVert}
\newcommand{\abs}[1]{\left\lvert#1\right\rvert}
\lstdefinestyle{engpy}{language=Python,basicstyle=\ttfamily\footnotesize,keywordstyle=\color{RoyalBlue}\bfseries,commentstyle=\color{gray},stringstyle=\color{teal},backgroundcolor=\color{pale},frame=single,rulecolor=\color{blue!30},breaklines=true,columns=fullflexible,keepspaces=true,showstringspaces=false,tabsize=4,upquote=true}
\lstset{style=engpy}
\begin{document}
\begin{titlepage}
\pagecolor{navy}\color{white}
\vspace*{0.75in}
{\sffamily\bfseries\large MATHEMATICS \& ENGINEERING DOCUMENTATION SUITE\\[12pt]}
{\color{cyan!40}\rule{\textwidth}{2pt}}\\[1.0in]
{\sffamily\bfseries\fontsize{29}{36}\selectfont Numerical Methods and Scientific Computing Guide\par}
\vspace{0.45in}
{\Large Volume 4 of 8}\par\vspace{0.55in}
{\color{cyan!30}\rule{0.50\textwidth}{1pt}}\par\vspace{0.3in}
{\large Independent technical companion aligned with the\par}
{\large 13-chapter organization of Nagle, Saff \& Snider (7th ed.)\par}
\vfill
{\large Original derivations, worked examples, verification and engineering practice\par}\vspace{0.3in}
{\small Prepared October 2026 \quad\textbullet\quad pdf\LaTeX{} compatible\par}
\end{titlepage}
\nopagecolor\color{ink}\pagenumbering{roman}
\chapter*{Reader's guide and source boundaries}\addcontentsline{toc}{chapter}{Reader's guide and source boundaries}
This independent study resource is organized around \emph{Fundamentals of Differential Equations and Boundary Value Problems}, Seventh Edition, by R.~Kent Nagle, Edward~B. Saff, and Arthur~David Snider (Pearson, 2018; ISBN 978-0-321-97710-6). The book's chapter and section numbers are given as \emph{navigation references}, not as claims of reproduced text or proofs.

\begin{bookground}
\textbf{Textbook-grounded scope:} topics, terminology, and chapter/section organization corresponding to the cited textbook. All explanations, derivations, exercise prompts, and diagrams in this volume are newly authored and may emphasize different examples. Check the textbook itself for its exact hypotheses, proofs, exercises, or nomenclature in a particular section.
\end{bookground}
\begin{modern}
\textbf{Distinct extension layer:} Python/NumPy/SciPy workflows, engineering QA and reproducibility practices, automated convergence studies, diagnostic plots, and implementation guidance. These are \emph{not} represented as content supplied by the textbook. API calls are examples and should be verified against the installed package version.
\end{modern}
\begin{caution}
\textbf{Mathematical qualification:} formulas require their stated regularity, domain and boundary/initial conditions. Numerical agreement does not replace a proof; computational outputs are not certified error bounds unless explicitly established.
\end{caution}
\tableofcontents\clearpage\pagenumbering{arabic}
\chapter{Numerical initial-value methods}
\begin{bookground}Correspondence: \S1.4 Euler approximation, \S3.6 closer look at Euler's algorithm, \S3.7 Taylor and Runge--Kutta methods, \S5.3 numerical systems, Appendix F on Runge--Kutta for $n$ equations.\end{bookground}
\section{Discrete approximation and error vocabulary}
Let $y'=f(t,y)$ with $y(t_0)=y_0$ on a grid $t_n=t_0+nh$. Numerical state $y_n$ approximates $y(t_n)$. \emph{Local truncation error} describes the per-step defect produced by inserting the exact solution into an update; \emph{global error} compares the propagated approximation with the exact solution. \emph{Consistency}, \emph{stability}, and \emph{convergence} are related but not interchangeable.
\section{Forward Euler and its limitations}
Euler's algorithm is $y_{n+1}=y_n+h f(t_n,y_n)$. For sufficiently smooth $f$ and solution, the method has global order one on fixed finite intervals under standard stability assumptions. The explicit update applied to $y'=\lambda y$ gives $y_{n+1}=(1+h\lambda)y_n$; decay of the numerical mode requires $|1+h\lambda|<1$ for nonzero initial data. Even for an analytically stable system, a large step can induce numerical growth.
\begin{example}[Euler error calculation]
Solve $y'=-y$, $y(0)=1$ over $t=1$ using $h=1/4$: $y_4=(3/4)^4=81/256\approx0.31640625$. The exact solution is $e^{-1}\approx0.36787944$. The absolute error is approximately $0.05147$.
\end{example}
\section{Classical fourth-order Runge--Kutta}
For $k_1=f(t_n,y_n)$, $k_2=f(t_n+h/2,y_n+hk_1/2)$, $k_3=f(t_n+h/2,y_n+hk_2/2)$, and $k_4=f(t_n+h,y_n+hk_3)$, update
\[y_{n+1}=y_n+\frac h6(k_1+2k_2+2k_3+k_4).\]
Under smoothness and stability assumptions the global order is four. RK4 is not automatically superior for stiff equations, discontinuities or singular solutions; error control requires either step refinement or embedded methods.
\section{Adaptive solvers and stiffness}
Modern adaptive ODE solvers estimate local error and change step size. Explicit Runge--Kutta methods suit many nonstiff systems. Implicit methods such as Radau and BDF are important when fast decaying modes constrain explicit stability despite much slower dynamics of interest. \emph{Stiffness} is a property of the model, time interval, chosen accuracy and algorithm; it is not established solely by nonlinearity.
\begin{modern}SciPy \texttt{solve\_ivp} supports explicit RK23/RK45/DOP853 and implicit Radau/BDF among others; choose absolute and relative tolerances with state scales in mind. The interface and library names are modern extensions, not Chapter 3 textbook statements.\end{modern}
\chapter{Implement and validate solvers}
\section{Reference implementation of Euler and RK4}
\begin{lstlisting}
import numpy as np

def euler(f, t0, y0, h, steps):
    t = float(t0)
    y = np.asarray(y0, dtype=float).copy()
    out = [(t, y.copy())]
    for _ in range(steps):
        y = y + h * np.asarray(f(t, y))
        t = t + h
        out.append((t, y.copy()))
    return out

def rk4(f, t0, y0, h, steps):
    t = float(t0)
    y = np.asarray(y0, dtype=float).copy()
    out = [(t, y.copy())]
    for _ in range(steps):
        k1 = np.asarray(f(t, y))
        k2 = np.asarray(f(t+h/2, y+h*k1/2))
        k3 = np.asarray(f(t+h/2, y+h*k2/2))
        k4 = np.asarray(f(t+h, y+h*k3))
        y = y + h * (k1+2*k2+2*k3+k4) / 6
        t = t + h
        out.append((t, y.copy()))
    return out
\end{lstlisting}
\section{Step-refinement convergence study}
For $h,h/2,h/4$, compute comparable endpoint errors $E_h$. An observed order estimate is $p_{\mathrm{obs}}=\log_2(E_h/E_{h/2})$ when the errors are nonzero and follow the asymptotic model $E_h\approx C h^p$. Without an exact solution, Richardson-style differences can approximate rates but require careful assumptions. Do not compute rate ratios from unrelated output points or tolerances.
\begin{lstlisting}
from math import exp, log2
for n in [10, 20, 40, 80]:
    result = rk4(lambda t, y: -y, 0.0, [1.0], 1.0/n, n)
    error = abs(result[-1][1][0] - exp(-1.0))
    print(n, error)
\end{lstlisting}
\section{A production numerical QA checklist}
\begin{enumerate}
\item Verify function dimensions and units, and record parameter provenance.
\item Check solver status, finite states, error estimates and termination reason.
\item Refine tolerances or mesh; compare an independent method when practical.
\item Test known equilibria, conservation laws, positivity and response symmetry.
\item Localize sharp transitions and events instead of relying on sparse samples.
\item Track solver and dependency versions; make random seeds explicit if present.
\end{enumerate}
\chapter{Coupled systems and nonlinear phenomena}
\begin{bookground}Correspondence: \S5.3 numerical systems, \S5.4 phase plane, \S5.8 Poincar\'e maps and chaos, and \S12.6 periodic solutions.\end{bookground}
\section{Rewrite higher order as first order}
For $y''+2\zeta\omega_0 y'+\omega_0^2 y=F(t)$, set $x_1=y$, $x_2=y'$. Then $x_1'=x_2$ and $x_2'=F(t)-2\zeta\omega_0x_2-\omega_0^2x_1$. Nearly all modern ODE solver interfaces accept this first-order vector form.
\section{Phase planes, events and invariants}
For autonomous $\mathbf x'=\mathbf f(\mathbf x)$ in two dimensions, plotted trajectories and a direction field provide qualitative insight but cannot prove stability without supporting analysis. For a mechanical system with energy $E=\tfrac12m v^2+\tfrac12 k x^2$, differentiate along its trajectories; undamped unforced motion has $\dot E=0$, while damping gives $\dot E=-cv^2\le0$.
\section{Events and transition times}
An event function $g(t,\mathbf x(t))=0$ may detect a threshold or zero crossing. Event detection generally relies on sign changes over solver steps; very rapid repeated crossings may be missed. Validate event accuracy under step-size restrictions and report whether detected events terminate integration.
\chapter{Boundary-value problems: shooting and collocation}
\begin{bookground}Correspondence: \S11.4--11.6 nonhomogeneous boundary-value problems and Green's functions; Chapter 11 projects D and E on shooting and finite differences. The SciPy collocation API is an independent modern extension.\end{bookground}
\section{Shooting method}
Given $y''=f(x,y,y')$, $y(a)=\alpha$, $y(b)=\beta$, choose an unknown slope $s=y'(a)$. Numerically propagate the IVP and define $R(s)=y_s(b)-\beta$. Solve $R(s)=0$ by a root method. Multiple or absent roots reflect the underlying BVP, not necessarily the root-finder implementation. Sensitive IVPs can make shooting poorly conditioned.
\begin{example}[Analytically soluble shooting benchmark]
For $y''=-2$, $y(0)=0$, $y(1)=0$, integration gives $y(x)=-x^2+s x$; $R(s)=s-1$, so $s=1$, $y=x(1-x)$. This is a convenient code-validation case.
\end{example}
\section{Collocation BVP solver}
A first-order vector BVP $\mathbf y'=\mathbf f(x,\mathbf y)$ has a residual form $\mathbf r(\mathbf y(a),\mathbf y(b))=0$. For $n$ state components and $k$ unknown parameters, a determined formulation requires $n+k$ independent boundary conditions. This counts constraints, but uniqueness still requires analysis.
\begin{modern}SciPy \texttt{solve\_bvp} takes a mesh, initial profile guess, differential-equation callback and boundary residual callback. The status code and boundary residual should be examined, and a refined mesh should be used for numerical verification.\end{modern}
\begin{lstlisting}
import numpy as np
from scipy.integrate import solve_bvp

# y'' = -2, y(0) = y(1) = 0; exact y = x*(1-x)
def ode(x, Y):
    return np.vstack((Y[1], -2*np.ones_like(x)))

def bc(Ya, Yb):
    return np.array([Ya[0], Yb[0]])

x = np.linspace(0.0, 1.0, 7)
Y = np.zeros((2, x.size))
sol = solve_bvp(ode, bc, x, Y, tol=1e-7)
assert sol.success, sol.message
xx = np.linspace(0, 1, 101)
assert max(abs(sol.sol(xx)[0] - xx*(1-xx))) < 1e-5
\end{lstlisting}
\chapter{Finite differences for PDEs and spatial BVPs}
\begin{bookground}Correspondence: Chapter 10 project D (numerical Poisson problem on a rectangle), Chapter 11 project E (finite difference for BVPs), and Chapter 10 heat/wave/Laplace equations.\end{bookground}
\section{Central differences}
At a smooth interior point $x_i$ with uniform step $h$,
\[y''(x_i)\approx\frac{y_{i-1}-2y_i+y_{i+1}}{h^2},\qquad y'(x_i)\approx\frac{y_{i+1}-y_{i-1}}{2h}.\]
Both have local truncation error $O(h^2)$ under sufficient smoothness. Applying the second derivative to $-y''=f$ on interior nodes gives a tridiagonal positive-definite matrix for zero Dirichlet endpoints.
\section{One-dimensional Poisson benchmark}
For $-y''=2$ on $(0,1)$ and $y(0)=y(1)=0$, the exact solution $y=x(1-x)$ is quadratic, so its second central difference is exact on any uniform grid in exact arithmetic. This tests assembly and boundary handling but cannot reveal the asymptotic order (an exact polynomial benchmark may mask truncation error).
\section{Method of lines for the heat equation}
For $u_t=\kappa u_{xx}$ on a fixed spatial grid, replace the second spatial derivative by a matrix $D_{xx}$. The semi-discrete ODE is $\dot{\mathbf u}=\kappa D_{xx}\mathbf u+\mathbf b(t)$. Explicit Euler has a familiar stability restriction $r=\kappa\Delta t/(\Delta x)^2\le1/2$ for the classical one-dimensional interior-centered scheme with Dirichlet conditions; this bound is not universal for arbitrary geometry, dimensions or time steppers.
\begin{caution}Stable does not mean accurate. Boundary handling, truncation error and the conditioning of spatial operators must also be checked.\end{caution}
\chapter{Sparse linear algebra and eigenvalue approximations}
\begin{modern}Sparse matrix methods and the \texttt{eigsh} interface are computational extensions that complement Chapter 9 matrix methods and Chapter 11 eigenvalue theory.\end{modern}
\section{Discrete Sturm--Liouville operator}
The Dirichlet eigenproblem $-y''=\lambda y$ on $(0,1)$ has exact $\lambda_n=n^2\pi^2$, $y_n=\sin(n\pi x)$. Central differences on $N$ interior points give the matrix $A=(1/h^2)\operatorname{tridiag}(-1,2,-1)$ with $h=1/(N+1)$. Its eigenvalues approximate the continuous eigenvalues from below for fixed mode index as $N\to\infty$.
\begin{lstlisting}
import numpy as np
from scipy.sparse import diags
from scipy.sparse.linalg import eigsh
N = 200
h = 1.0 / (N + 1)
A = diags([-np.ones(N-1), 2*np.ones(N), -np.ones(N-1)],
          [-1, 0, 1], format="csr") / h**2
lam, vec = eigsh(A, k=3, which="SM")
lam = np.sort(lam)
print("discrete:", lam)
print("continuous:", (np.arange(1,4)*np.pi)**2)
\end{lstlisting}
The eigenvector sign is arbitrary, so compare modes up to sign and normalization. For generalized eigenproblems use a consistent mass matrix and weighted orthogonality.
\chapter{Computational exercises and laboratory acceptance tests}
\section{Laboratory sequence}
\begin{enumerate}
\item Implement Euler and RK4 for $y'=-y$, compare errors at $h=0.1,0.05,0.025$; estimate observed order.
\item Simulate an underdamped oscillator; verify decreasing mechanical energy and parameter sensitivity.
\item Solve the BVP $y''=-2$, $y(0)=y(1)=0$ by shooting and collocation; compare residuals.
\item Approximate the first three eigenvalues of $-y''=\lambda y$ on $[0,1]$; halve $h$ and quantify the discrepancy.
\item Integrate $y'=-100(y-\cos t)-\sin t$, $y(0)=1$ with explicit and implicit methods; compare effort and errors against $y=\cos t$.
\end{enumerate}
\section{Expected invariants and rubric}
A successful report includes equations with units; initial/boundary conditions; analytic benchmark if available; solver settings and status; convergence table; error metric; plots with labeled axes; and a reasoned conclusion about numerical vs model error. For the stiff scalar benchmark in exercise 5, substitution shows $y=\cos t$ is exact, so any discrepancy is numerical. Compare solver work only under matched accuracy requirements.

\appendix
\chapter{References and reproducibility}
\begin{enumerate}
\item R.~K. Nagle, E.~B. Saff, and A.~D. Snider, \emph{Fundamentals of Differential Equations and Boundary Value Problems}, 7th ed., Pearson, 2018, ISBN 978-0-321-97710-6. Primary source for the \emph{topic map and terminology}; cited locally by chapter and section.
\item SciPy project, \emph{solve\_ivp} reference, \url{https://docs.scipy.org/doc/scipy/reference/generated/scipy.integrate.solve_ivp.html} (modern extension).
\item SciPy project, \emph{solve\_bvp} reference, \url{https://docs.scipy.org/doc/scipy/reference/generated/scipy.integrate.solve_bvp.html} (modern extension).
\item SciPy project, \emph{eigsh} reference, \url{https://docs.scipy.org/doc/scipy/reference/generated/scipy.sparse.linalg.eigsh.html} (modern extension).
\item Python, NumPy, SciPy, and SymPy project documentation; computational checks require recording package versions and tolerances.
\end{enumerate}
\section*{Verification record template}
\begin{tabularx}{\textwidth}{>{\bfseries}p{0.30\textwidth}X}\toprule
Problem identifier & \rule{0pt}{15pt}\ \\ \midrule
ODE/PDE and domain & \\[8pt]
Initial/boundary conditions & \\[8pt]
Units and parameter values & \\[8pt]
Analytic residual & \\ [8pt]
Grid/step refinement & \\ [8pt]
Independent solver comparison & \\ [8pt]
Version and reproducibility notes & \\ [8pt]\bottomrule
\end{tabularx}
\end{document}
